\begin{exercice}\label{exo5-1} (\minsyndical)

On a g\'en\'er\'e artificiellement un \'echantillon de $n=14$ individus en suivant la loi normale
d'esp\'erance $m=10$ et de variance $\sigma^2=4$. On donne ici les valeurs obtenues :

\[
12.1\quad9.8\quad10.8\quad11.9\quad5.8\quad8.7\quad8.0\quad9.6\quad13.0\quad9.9\quad12.5\quad8.6\quad10.0\quad8.4. 
\]

\begin{enumerate}

\item Calculer la moyenne empirique $\bar{x}_n$ et comparer \`a la valeur $m$ de l'esp\'erance.

\item Calculer la variance empirique $s^2_n$ et comparer \`a la valeur $\sigma^2$ de la variance.\\
Que remarque-t'on ?

\item Calculer la variance empirique corrig\'ee $\frac{n}{n-1} s^2_n$. Conclure.\\

\end{enumerate}

%\corrref{TD1_1}

\end{exercice}
